% Ch17 WS3 -- temperature dependence, dG and K, kinetic control. Blocks 3-4.
\documentclass{apchem-worksheet}

\apcunitword{Chapter}
\apcunit{17}
\apcunittitle{Entropy and Free Energy}
\apcdoctitle{Worksheet 3 \textbullet\ Temperature, $K$, and Kinetic Control}
\apcsubtitle{Zumdahl \S17.3, 17.9 \textbullet\ favored says where it ends up, never how fast}

\begin{document}
\wsheader[$\Delta G^\circ = \Delta H^\circ - T\Delta S^\circ$ \textbullet\
$\Delta G^\circ = -RT\ln K$ \textbullet\ $R = \SI{8.314}{\joule\per\mole\per\kelvin}$
\textbullet\ $RT \approx \SI{2.5}{\kilo\joule\per\mole}$ at
\SI{298}{\kelvin}. Remember to convert $\Delta G^\circ$ to joules before
using $R$.]

\problem[]{Complete the four-case table:}
\begin{parts}
  \item $\Delta H^\circ < 0$, $\Delta S^\circ > 0$: favored at
        \answerblank[5.4cm]{all temperatures}
  \item $\Delta H^\circ > 0$, $\Delta S^\circ < 0$: favored at
        \answerblank[5.4cm]{no temperature}
  \item $\Delta H^\circ > 0$, $\Delta S^\circ > 0$: favored at
        \answerblank[5.4cm]{high temperature}
  \item $\Delta H^\circ < 0$, $\Delta S^\circ < 0$: favored at
        \answerblank[5.4cm]{low temperature}
  \item In which two cases is no calculation needed at all, and why?
        \answerlines[2]{(a) and (b). When both terms point the same way the
        sign of $\Delta G^\circ$ is fixed regardless of $T$, so there is
        nothing to compute.}
\end{parts}

\problem[]{The crossover temperature is where $\Delta G^\circ = 0$.}
\begin{parts}
  \item Derive the expression for it from
        $\Delta G^\circ = \Delta H^\circ - T\Delta S^\circ$.
        \answerlines[2]{Setting $\Delta G^\circ = 0$ gives
        $0 = \Delta H^\circ - T\Delta S^\circ$, so
        $T = \Delta H^\circ/\Delta S^\circ$.}
  \item A reaction has $\Delta H^\circ = +30$~kJ and
        $\Delta S^\circ = +80$~J/K. Find the crossover temperature and say
        on which side it is favored. \workspace[2.2cm]
        \keyonly{\begin{apcanswer}$T = 30/0.080 = \mathbf{375}$~K
        ($\approx\SI{102}{\celsius}$). Both terms positive, so it is
        favored \textbf{above} 375 K.\end{apcanswer}}
  \item A reaction has $\Delta H^\circ = -40$~kJ and
        $\Delta S^\circ = -120$~J/K. Find the crossover temperature and
        the favored side. \workspace[2.2cm]
        \keyonly{\begin{apcanswer}$T = -40/-0.120 = \mathbf{333}$~K
        ($\approx\SI{60}{\celsius}$). Both terms negative, so it is
        favored \textbf{below} 333 K.\end{apcanswer}}
\end{parts}

\problem[]{Vaporization of water has $\Delta H^\circ = +44$~kJ and
$\Delta S^\circ = +119$~J/K.}
\begin{parts}
  \item Calculate $\Delta G^\circ$ at \SI{298}{\kelvin} and interpret it.
        \workspace[2.2cm]
        \keyonly{\begin{apcanswer}$44 - 298(0.119) = 44 - 35.5 =
        \mathbf{+8.5}$~kJ --- not favored, which is why a glass of water
        does not all become vapour at room
        temperature.\end{apcanswer}}
  \item Calculate the crossover temperature and identify what it physically
        is. \workspace[2.2cm]
        \keyonly{\begin{apcanswer}$T = 44/0.119 = \mathbf{370}$~K
        $\approx\SI{97}{\celsius}$ --- the \textbf{normal boiling
        point}\end{apcanswer}}
  \item Water actually boils at \SI{373}{\kelvin}. Account for the
        discrepancy. \answerlines[2]{The tabulated $\Delta H_f^\circ$ and
        $S^\circ$ values are rounded to the nearest whole unit, and the
        crossover is a ratio of two such rounded numbers. A 3 K gap on a
        370 K result is well within that rounding.}
\end{parts}

\problem[]{Convert between $\Delta G^\circ$ and $K$ at \SI{298}{\kelvin}:}
\begin{parts}
  \item $\Delta G^\circ = -20$~kJ. Find $K$. \workspace[2cm]
        \keyonly{\begin{apcanswer}$\ln K = 20{,}000/[(8.314)(298)] = 8.07$;
        $K = e^{8.07} = \mathbf{3.2\times10^{3}}$\end{apcanswer}}
  \item $\Delta G^\circ = +20$~kJ. Find $K$.
        \answerblank[3.4cm]{$3.1\times10^{-4}$}
  \item $K = 25$. Find $\Delta G^\circ$. \workspace[2cm]
        \keyonly{\begin{apcanswer}$\Delta G^\circ = -RT\ln K =
        -(8.314)(298)\ln 25 = -(2477)(3.22) = -7977~\text{J} =
        \mathbf{-8.0}~\text{kJ}$\end{apcanswer}}
  \item $K = 1.0\times10^{-5}$. Find $\Delta G^\circ$.
        \answerblank[3cm]{$+28.5$ kJ}
\end{parts}

\problem[]{Estimation without a calculator. Fill in ``greater than 1,''
``close to 1,'' or ``less than 1'':}
\begin{parts}
  \item $\Delta G^\circ = -45$~kJ: $K$ is
        \answerblank[4.4cm]{much greater than 1}
  \item $\Delta G^\circ = +45$~kJ: $K$ is
        \answerblank[4.4cm]{much less than 1}
  \item $\Delta G^\circ \approx 0$: $K$ is
        \answerblank[4.4cm]{close to 1}
  \item Why does a change of only a few kJ in $\Delta G^\circ$ swing $K$ so
        violently? \answerlines[3]{$\Delta G^\circ$ enters through an
        exponential, $K = e^{-\Delta G^\circ/RT}$, and the natural scale it
        is compared against is $RT \approx \SI{2.5}{\kilo\joule\per\mole}$
        at \SI{298}{\kelvin}. A change of a few kJ is therefore a change of
        one or two in the exponent, which multiplies or divides $K$ several
        times over.}
\end{parts}

\problem[]{For \ce{N2O4(g) -> 2NO2(g)}, $\Delta H^\circ = +58$~kJ and
$\Delta S^\circ = +176$~J/K.}
\begin{parts}
  \item Calculate $\Delta G^\circ$ at \SI{298}{\kelvin}.
        \answerblank[3cm]{$+5.6$ kJ}
  \item Calculate $K$ at \SI{298}{\kelvin}. \workspace[2.2cm]
        \keyonly{\begin{apcanswer}$\ln K = -5600/[(8.314)(298)] = -2.26$;
        $K = \mathbf{0.10}$ --- reactant-favored, but only
        modestly\end{apcanswer}}
  \item Find the crossover temperature and explain what a student would
        observe on warming a sealed tube of this mixture.
        \answerlines[3]{$T = 58/0.176 = 330$~K ($\approx\SI{57}{\celsius}$).
        Above that, dissociation becomes favored, so warming the tube
        shifts the equilibrium towards \ce{NO2} and the mixture darkens ---
        \ce{NO2} is brown, \ce{N2O4} colourless.}
\end{parts}

\problem[]{Kinetic control.}
\begin{parts}
  \item Define kinetic control and give the usual cause.
        \answerlines[2]{A process that is thermodynamically favored but
        proceeds too slowly to observe is under kinetic control. The usual
        cause is a high activation energy.}
  \item The conversion of diamond to graphite has
        $\Delta G^\circ < 0$ at \SI{298}{\kelvin}. Explain why diamonds
        persist. \answerlines[3]{The reaction is favored but requires
        breaking and reorganizing the entire covalent network, so the
        activation energy is enormous. At room temperature effectively no
        collisions supply it, and the rate is immeasurably slow. It is
        under kinetic control, not at equilibrium.}
  \item A mixture of \ce{H2} and \ce{O2} sits unchanged for years, yet the
        combustion has $\Delta G^\circ = -474$~kJ. What does a spark change,
        and what does it not? \answerlines[3]{The spark supplies energy to
        surmount the activation barrier locally, raising the rate. It does
        not change $\Delta H^\circ$, $\Delta S^\circ$, $\Delta G^\circ$ or
        $K$ --- the reaction was always favored; it was only slow.}
  \item A catalyst is added to a favored but slow reaction. State the
        effect on the rate and on $K$.
        \answerblank[6.4cm]{rate increases; $K$ unchanged}
\end{parts}

\problem[]{A student writes: ``The reaction has a large negative
$\Delta G^\circ$, so it must go quickly and go to completion.'' Identify
both errors. \answerlines[4]{First, $\Delta G^\circ$ says nothing about
rate --- a strongly favored reaction can be under kinetic control and take
geological time. Second, ``favored'' does not mean ``complete'': it means
$K > 1$, so products dominate at equilibrium, but some reactant always
remains. Only an infinite $K$ would correspond to going to completion.}}

\begin{spiralstrip}{Chapter 17 \textbullet\ blocks 1--2}
  \item $\Delta S^\circ_{\text{rxn}} = $
        \answerblank[4.6cm]{$\sum S^\circ_{\text{prod}} -
        \sum S^\circ_{\text{react}}$}
  \item $S^\circ$ for an element is: \answerblank[2cm]{not zero}
  \item Sign of $\Delta S^\circ$ when gas moles rise:
        \answerblank[1.8cm]{positive}
  \item Units of $S^\circ$: \answerblank[3.4cm]{J/(K$\cdot$mol)}
  \item $\Delta S_{\text{surr}}$ matters most at what temperature?
        \answerblank[2cm]{low}
\end{spiralstrip}

\end{document}
